\chapter{讨论}

前一章以四幕结构呈现了实验结果的定量证据。本章从机制解释、方法论意义和研究局限性三个层面对这些结果展开讨论，目的是回答"为什么会出现这些现象"以及"这些发现的适用边界在哪里"。

\section{捷径竞争机制的解释}

S1阶段的公式坍缩现象和S2阻断实验的$\Delta\VER$结果共同支持如下解释：在当前实验设置下，包含自回归滞后特征的KAN符号提取更容易因为滞后项在稀疏化竞争中形成信息捷径，而系统性压制物理气象变量；这比"变量与预测目标无关"更能解释本文观察到的现象。

\textbf{自相关结构下的方差支配。}在5分钟采样分辨率下，净负荷序列呈现极强的短程自相关性。以load滞后序列为代表的自回归输入与当前净负荷保持高相关，其中更近的滞后项通常解释了短期预测目标方差的绝大部分。这意味着在任何以均方误差为损失函数的训练框架中，与滞后负荷相关的网络边只需以近乎恒等映射的方式传递信号，即可实现极低的训练损失。相比之下，气象物理变量对净负荷变化的边际贡献远小于滞后项的直接贡献——尽管从物理因果链条上看，正是这些气象因素驱动了风光出力的变化。

\textbf{L1+熵正则化下的竞争压力。}KAN的稀疏化训练阶段施加了L1范数和熵正则化（$\lambda_{L1}=1.0$，$\lambda_{\text{entropy}}=2.0$）。两者的联合效应本质上是在有限的"活跃边预算"下进行特征竞争：模型被迫将有限的表达能力分配给对损失函数降低贡献最大的特征通路。在这一竞争中，解释了90\%以上方差的滞后项必然胜出，而贡献剩余不到10\%方差的气象变量则在迭代优化过程中被逐步压制至剪枝阈值以下。这不是KAN架构本身的设计缺陷，而是任何具有稀疏化正则的参数化模型在面对方差贡献高度不平衡的特征集时都会表现出的结构性行为。

\textbf{特征遮蔽效应的类比。}上述机制在传统特征选择文献中有对应物。在Lasso回归和弹性网络中，当两个高度共线的特征同时进入模型时，正则化倾向于保留其中一个而将另一个的系数压缩至零——即所谓的"遮蔽效应"（masking effect）。本研究中，自回归滞后项不仅与目标高度相关，而且通过线性传递间接包含了气象因素的影响（因为历史净负荷本身就是气象条件作用的结果）。因此，滞后项实质上扮演了气象信息的"充分代理"角色。

\textbf{S2阻断实验提供的干预性证据。}Case 3（聚焦风电任务）中$\Delta\VER=+0.60$，95\% bootstrap置信区间$[0.20, 1.00]$不含零；这是本文最强的干预证据。相比之下，严格matched的Case 4（直接净负荷任务）只在4个有效paired seed上得到$\Delta\VER=+0.25$，置信区间为$[0.00, 0.75]$，下界触零；4个有效pair中仅seed5恢复了1条wind\_speed\_10m\_m\_s\_cubed边，另有seed1的unblocked run在两次CPU尝试中均未产出post-prune checkpoint。因此，Case 4现在只能提供弱且不稳定的direct-task支持，而不再被写成强干预确认。

\textbf{机制的普遍性。}值得强调的是，自回归捷径竞争并非KAN所特有的问题，而是高频时间序列预测中的一般性结构特征。任何将自回归滞后与外生变量共同输入的稀疏模型——无论是Lasso、稀疏MLP还是KAN——都面临同样的竞争压力。KAN的特殊之处在于其稀疏化-剪枝-符号拟合的三阶段提取管线使这一竞争的后果以"公式坍缩"的形式可视化呈现。

\section{实验过程的收敛链条}

本文的主结论并不是一开始就直接得到的。项目早期首先比较KAN与PySR等方法的预测性能，但需要区分历史探索期与当前主实验期：历史探索期在个别弱配置、非canonical的load目标run中曾出现PySR优于KAN的现象；进入当前canonical matched $\Delta\text{net\_load}_{h6}$主实验设置后，KAN teacher在预测层面已稳定优于direct PySR。因此，研究的关键问题从"哪种方法精度更高"转向"高精度teacher能否转化成保留物理因子的符号公式"。

direct路线的首个异常并不是wind，而是整体公式迅速退化为负荷主导。这一现象最初并不能自动推出"捷径竞争"解释，因为还存在另一种朴素可能：超参数、剪枝阈值或符号库设置得不合适，导致物理量偶然没有进入公式。正因如此，后续才需要focused teacher与component teacher式的定向实验，去验证"是不是所有物理量都进不去公式"。solar、GHI和temperature在多条定向路线中能够稳定进入局部公式，说明当前链路并非整体失效。

真正持续异常的是wind。即使在多轮调参、特征重组和teacher设计之后，wind仍更常表现为：原始风速项未显式进入公式，或只通过wind\_lag\_*、代理变量和替代表达间接出现。更稳妥的解释不是"风电任务训练失败"，而是wind在当前任务和时间尺度下具有更低的符号可辨识性。进一步的long-horizon结果也更适合写成wind可辨识性的非单调变化，即中等horizon更容易显现结构、整体上接近"先增强、后减弱"；至于其他物理量，现有证据只支持"更长horizon下整体更不稳定、部分量减弱"，不支持"持续下降"或"严格单调下降"的强表述。

这一收敛链条也界定了本文的叙事边界。本文讨论的是``KAN teacher + symbolic extraction''的后验提取流程，而不是经典教师-学生蒸馏；S3已经完成净负荷总公式的显式合成与test复算，但结果显示总公式对象与结构化组合预测器对象之间仍有明显transfer gap。因而，本文的主贡献应落在direct路线为何失败、S2如何提供干预性支持、以及S3如何恢复局部可辨识性并暴露总公式迁移间隙，而不是把"公式是否存在"误写成尚未完成的工程问题。

\section{精度与可辨识性的权衡}

S2阻断实验在揭示捷径竞争机制的同时，也暴露了一个不可回避的工程权衡：消除自回归捷径恢复了物理变量的可辨识性，但同时造成了预测精度的显著下降。

\textbf{精度代价的量化。}Case 4的blocked run清楚展示了这一权衡：直接阻断会带来明显精度代价。由于该RMSE来自delta域剪枝网络评估，而主表RMSE来自abs(test)教师模型评估，本文不再报告跨表百分比退化。

\textbf{聚焦任务中的反常现象。}Case 3（聚焦风电任务）的结果呈现了一个值得关注的反常现象：阻断自回归滞后后，模型RMSE反而从1363下降到1047。这一结果表明，在窄聚焦的单变量预测任务中，自回归滞后特征可能引入噪声而非信号。

\textbf{S3分解作为折中方案。}S3结构化组合预测器的技能分数skill=0.515，已经优于直接KAN教师（skill=0.453），RMSE为1254.62；同时三个子任务的FAR均达到3/3。S3方案之所以能同时保持精度和可解释性，关键在于子任务分解从结构层面缓解了竞争压力。需要再次强调的是，S3现在同时包含"三条局部公式 + 结构化组合预测器 + 已复算的净负荷总公式"三类对象：组合预测器保持了精度，而显式总公式在test复算中的RMSE为2644.36、skill为-0.021，暴露出相对预测器的明显transfer gap。

\textbf{对实际应用的启示。}如果运营目标是最大化短期预测精度（如实时调度），直接包含自回归特征的模型仍是最优选择；如果运营目标是理解预测背后的物理驱动因素（如中长期规划、异常诊断），S3分解方案提供了在可接受精度范围内获取物理可解释公式的可行路径。两种范式并非互斥，而是面向不同运营需求的互补工具。

\section{与MCKAN的方法论对应}

本研究的实验设计在多个层面上与Chen等\cite{chen2025mckan}提出的MCKAN存在方法论上的对应关系。

\textbf{多步预测与多时域差分目标。}MCKAN采用多步预测架构以覆盖15至45分钟的预测时域。本研究同样设计了多时域评估（$h=6, 12, 24$），但目标不在于各时域的精度排名，而在于观察可辨识性指标随预测时域的变化模式。

\textbf{特征选择与可辨识性诊断。}MCKAN使用Pearson相关系数作为特征选择依据。本研究中定义的VER和FAR在概念上构成了一种"深层特征选择"诊断——它们衡量的不是变量与目标的统计相关性，而是变量在经历完整的KAN稀疏化-剪枝-符号拟合管线后是否仍然存活。

\textbf{子任务分解结构。}MCKAN将风电和光伏预测作为独立子任务处理。本研究的S3结构化分解与之一脉相承，关键的方法论推进在于S3分解的设计动机不是提升精度，而是利用物理上自然的任务边界来消解自回归竞争压力。

\textbf{消融实验与阻断干预。}MCKAN采用模块消融实验来验证各模块的贡献。本研究的S2阻断实验在方法论上与之对应，但干预对象从网络模块转移到输入特征空间。

综合以上四个维度的对应关系，可以将本研究定位为MCKAN研究路线在可解释性维度上的延伸：MCKAN回答了"KAN能否高精度预测风光出力"的问题，本研究则尝试回答"KAN提取的符号公式能否保留物理含义"这一后续问题。

\section{局限性}

本研究在实验设计、数据条件和方法范围上存在以下局限性。

\textbf{第一，单一数据集的泛化性限制。}本研究全部实验基于ARPA-E PERFORM数据集的ERCOT区域数据。自回归捷径竞争机制作为结构性现象，理论上应在其他高频时间序列数据中同样存在，但其严重程度可能因地理区域、气候带、电网构成和采样分辨率的不同而有所差异。

\textbf{第二，种子数量的统计功效。}S2阻断实验使用了5个随机种子。虽然bootstrap置信区间在当前种子数下已排除了零效应，但区间宽度仍然较大。更大的种子数量将提供对$\Delta\VER$效应量更精确的估计。

\textbf{第三，数据预处理中的插值泄漏风险。}线性插值在训练-验证-测试划分之前执行，构成轻微的信息泄漏。然而，这一泄漏是非变量特异性的，不影响S2阻断实验中$\Delta\VER$的内部比较有效性；相应地，本文对"主要机制"的表述也保持在当前证据强度内，不作过度泛化。

\textbf{第四，缺乏物理约束的训练机制。}本研究中KAN的训练完全基于数据驱动，未融入任何物理先验知识。物理信息KAN（PIKAN）有望在不牺牲数据拟合能力的前提下提高符号提取的物理合理性，但这一方向超出了本研究的范围。

\textbf{第五，公式复杂度与人类可读性。}虽然S3方案成功使物理变量进入了提取公式，但所得公式的复杂度仍然较高，与"人类一眼可读"的简洁公式仍有差距。

\textbf{第六，KAN与PySR比较的公平性。}两者的优化空间存在本质差异：KAN通过梯度下降优化连续的B样条参数，而PySR通过遗传编程在离散的表达式空间中搜索。因此skill分数的直接比较不是严格的算力匹配比较。

上述局限性为后续研究指出了明确的改进方向，将在下一章中予以具体讨论。
